% Cas d'utilisation « Soumettre un plan CSR pour validation » — inclus sous la sous-section 4.4.3
\begin{table}[H]
  \centering
  \small
  \caption{Description textuelle du cas d'utilisation : «~Soumettre un plan CSR pour validation~»}
  \label{tab:cu-soumettre-plan-csr}
  \PFETableSetup
  \PFETableRowColors
  \PFETableArrayStretch
  \PFETableTabularXBegin
  \begin{tabularx}{\PFETableBlockWidth}{@{}>{\raggedright\arraybackslash}p{3.2cm} >{\raggedright\arraybackslash}X@{}}
    \tableHeaderStyle \tableHeaderCell{Champ} & \tableHeaderCell{Contenu} \\
    \textbf{Titre} & Soumettre un plan CSR pour validation \\
    \textbf{Acteurs} & Utilisateur Corporate ; utilisateur Site. \\
    \textbf{Objectif} &
      Permettre à l'utilisateur de soumettre un plan CSR afin de lancer le processus de validation et de figer les informations avant son exécution. \\
    \textbf{Préconditions} &
      1. L'utilisateur est authentifié.\par
      2. Un token est obtenu pour accéder aux ressources protégées.\par
      3. Un plan CSR existe.\par
      4. Le plan est complet et en état modifiable. \\
    \textbf{Postconditions} &
      1. Le plan CSR passe à l'état «~soumis~».\par
      2. Il devient non modifiable selon les règles métier.\par
      3. Il est prêt pour le processus de validation.\par
      4. Une notification est envoyée à l'utilisateur corporate. \\
    \textbf{Scénario nominal} &
      1. L'utilisateur accède au module des plans CSR.\par
      2. Le système affiche la liste des plans existants.\par
      3. L'utilisateur sélectionne un plan.\par
      4. Le système affiche les détails du plan.\par
      5. L'utilisateur déclenche l'action de soumission.\par
      6. Le système vérifie que les informations requises sont complètes.\par
      7. Le système change le statut du plan en «~soumis~».\par
      8. Le système enregistre la modification. \\
    \textbf{Scénario alternatif} &
      \textbf{A. Données incomplètes :} \par
      A.1 Le système refuse la soumission. \par
      A.2 Il signale les champs ou règles non satisfaits. \par\medskip
      \textbf{B. Plan non modifiable :} \par
      B.1 Le système informe l'utilisateur que la soumission n'est pas possible. \par
      B.2 Le message indique la cause (statut, droits ou fenêtre de modification). \\
  \end{tabularx}%
  \PFETableTabularXEnd
\end{table}
